% ---------------------------------------------------------------------------
% Western Digital - Firmware Engineering Intern (umbrella early-career req)
% Tailored from bases/embedded_controls.tex on 2026-08-03
% Worksheet: notes.md in this folder. Posting: job_description.txt
%
% Light touch on purpose - this is a pipeline req covering five tracks, so it
% does not earn deep tailoring. Fixed-point PI promoted to lead; assembly kept
% prominent (they name it); STM32H7/SPI dashboard kept as the clearest
% firmware-touching-hardware evidence.
%
% Do NOT put the company name anywhere in the rendered document. It belongs in
% the filename only.
% ---------------------------------------------------------------------------
\documentclass[10pt,letterpaper]{article}
\usepackage{resume}

\begin{document}
\bodyfont

\resumeheader
\educationsection

\sectionheader{TECHNICAL SKILLS}

\techline{\textbf{Languages}: C, C++, Assembly, Python, Java}
\techline{\textbf{Embedded \& Firmware}: TTC 60 VCU, STM32H7, PX4, real-time firmware, fixed-point arithmetic, deterministic execution, CAN, SPI, microcontroller I/O}
\techline{\textbf{Controls \& Algorithms}: Fixed-point PI/PID, anti-windup, saturation, slew-rate limiting, state estimation, sensor fusion}
\techline{\textbf{Test \& Tools}: Software-in-the-loop replay, regression tests, pytest, GDB, TASKING, TTC Downloader, BusMaster, MoTeC i2 Pro, Git}

\sectionheader{EXPERIENCE}

\roleline{Software Lead}{Jun 2026 - Present}
\orgline{\spartanracing}
\begin{itemize}
  \item Lead software planning for VCU controls, embedded firmware, telemetry, CAN diagnostics, and validation workflows
  \item Coordinate priorities and technical reviews across power limiting, launch control, dashboard, simulation, and vehicle testing
\end{itemize}

\entryspace

\roleline{Software Engineer Intern}{Feb 2026 - Present}
\orgline{Argus Defense \textbar{} San Jose, CA}
\begin{itemize}
  \item Developed and validated embedded flight-control firmware in C on PX4, iterating control gains against logged flight telemetry
  \item Integrated and tuned visual-inertial odometry for onboard state estimation, improving position-hold accuracy
  \item Tuned PID flight-control loops to reduce vertical and horizontal position drift in autonomous flight
\end{itemize}

\entryspace

\roleline{Software Controls Engineer}{Jul 2025 - Jun 2026}
\orgline{\spartanracing}
\begin{itemize}
  \item Implemented a torque-to-power controller in C with fixed-point PI control, saturation, anti-windup, and slew limits for deterministic TTC 60 VCU execution
  \item Built CAN diagnostics firmware for live power and torque frames, status bits, and error counters, debugging bus behavior against real-time BusMaster captures
  \item Designed a lap-adaptive energy-management algorithm setting a power setpoint from cumulative energy versus budget, holding actual power within \textbf{0.01\%} of setpoint
\end{itemize}

\entryspace

\roleline{Software Intern}{Aug 2024 - Jun 2025}
\orgline{\spartanracing}
\begin{itemize}
  \item Built deterministic entry and exit state logic for power limiting from measured power, requested torque, and threshold flags
  \item Tuned the PI torque controller used for power limiting to within \textbf{2\%} error, smoothing torque transitions to preserve drivability
  \item \textbf{Results:} 1st in endurance at MIS EV 2025
\end{itemize}

\sectionheader{PROJECT EXPERIENCE}

\roleline{Formula SAE EV Dashboard}{Feb 2026 - Present}
\orgline{\spartanracing}
\techline{C, STM32H7, CAN, SPI}
\begin{itemize}
  \item Built a real-time CAN telemetry display on an STM32H7 driving an LCD over SPI, decoding HV voltage, throttle, cell temperature, and power budget
  \item Designed an energy bar and LED alert system giving the driver lap-by-lap feedback on energy budget
\end{itemize}

\entryspace

\roleline{SR-Wireless-CAN --- Wireless Firmware Flashing \& Telemetry}{2026}
\orgline{\spartanracing \textbar{} Backend lead}
\techline{Python, python-can, CAN trace analysis, FastAPI, WebSockets, SQLite, Docker}
\begin{itemize}
  \item Proposed and implemented wireless firmware flashing for the TTC 60 VCU, decoding the undocumented flashing sequence from CAN trace captures and reimplementing it in Python over the bus
  \item Led backend development for the telemetry platform, streaming live CAN signals to a browser dashboard with DBC-driven decoding and recorded-trace replay
\end{itemize}

\end{document}
